\scanpage{65}
\mspage{9}
% archive sheet number: 64
Le troisième est un foncteur canonique
\[
\unsure{e|c|C}:\mathrm{Cat}_{/S}\longrightarrow\underline{\unsure{\mathrm{Parf}|\mathrm{Perf}}}\,S,
\]
qu'on va expliciter également.

Le foncteur $\Phi$ (resp. $\Psi$) appliqué à $\underline{\mathrm{Fib}}_0\,S$
(resp. $\mathrm{Cofib}_0\,S$), \deleted{\unsure{donne}} \inserted{induit} un foncteur
\[
\underline{\mathrm{Fib}}_0\,S\longrightarrow\underline{\mathrm{Fibsc}}_0
\qquad\text{respt.}\quad
\underline{\mathrm{Cofib}}_0\,S\longrightarrow\underline{\mathrm{Cofibsc}}_0\,S
\tag*{$(*)$}
\]
(ce \uline{qu'il faudra vérifier}). Cela fournit toutes
les flèches ``inverses'' dont on a besoin, à l'exception seulement des $\alpha,\beta,\alpha',\beta'$. Mais si on a
prouvé, \unsure{par} la méthode indiquée, que toutes les
flèches du diagramme induisent des équivalences sur les catégories localisées, pour prouver
qu'il en est encore de $=$ pour $\deleted{\unsure{\alpha}},\beta,\alpha',\beta'$, il
suffit de le prouver pour les \inserted{deux} foncteurs :
\[
\underline{\mathrm{Fibsc}}_0\,\deleted{(}S\deleted{)}\longrightarrow\underline{\mathrm{Fib}}\,S,\qquad
\underline{\mathrm{Cofibsc}}_0\,S\longrightarrow\underline{\mathrm{Cofib}}\,S\,;
\]
\unsure{or|où} il reste à définir des foncteurs en sens
inverse
\[
\underline{\mathrm{Fib}}\,\deleted{(}S\deleted{)}\to\underline{\mathrm{Fibsc}}_0\,\deleted{(}S\deleted{)},\qquad
\underline{\mathrm{Cofib}}\,S\unsure{\text{---}|\to}\mathrm{Cofibsc}_0\,S,
\]
et comme les sources sont des sous-catégories
pleines de $\mathrm{Cat}_{/S}$, il suffira de définir
\[
\Phi_0:\mathrm{Cat}_{/S}\longrightarrow\underline{\mathrm{Fibsc}}_0\,S,\qquad
\Psi_0:\mathrm{Cat}_{/S}\longrightarrow\mathrm{Cofib}\deleted{\unsure{s}}\mathrm{sc}_0\,S.
\]
Pour ceci, on posera
\[
\Phi_0=\Phi\unsure{\circ|*}C,\qquad\Psi_0=\Psi\circ C,\tag*{$\unsure{(**)|(xx)}$}
\]
